\cotefile{123}
\scanpage{23}
are isomorphic to subgroups of the group of exterior automorphisms of $\mathrm{G}$) ?

In a different direction, it would be\handwritten{\raisebox{.6ex}{$\vert$}}interesting to know more about the canonical morphism $\underline{\mathrm{g}}\handwritten{\longrightarrow}\underline{\mathrm{h}}/\mathrm{W}$, where $\underline{\mathrm{g}}$ is the Lie algebra of $\mathrm{G}$ and $\underline{\mathrm{h}}$ a Cartan subalgebra, and of the orbits of $\mathrm{G}$ on $\underline{\mathrm{g}}$, and the analoguous morphism of Steinberg $\mathrm{G}\handwritten{\longrightarrow}\mathrm{T}/\mathrm{W}$ and the orbits of $\mathrm{G}$ acting on itself. Do you know for instance exactly how varies the rank (= rang of the tangent map) of this \deleted{\unsure{applicati}} map ? There are some reasons to believe that \deleted{this} for points of unipotent orbits of dimension $\inserted{\mathrm{n}-}\mathrm{r}-2$ (just \deleted{the next lower after}\inserted{the next lower after} the maximal dimension $\mathrm{n}-\mathrm{r}$), the rank is $\mathrm{r}-1$ in case all roots have same length (that is the corresponding diagrams $\mathrm{A}_{\deleted{\unsure{\mathrm{n}|\mathrm{r}}}\unsure{\mathrm{r}|\mathrm{n}}},\mathrm{D}_{\deleted{\unsure{\mathrm{n}|\mathrm{r}}}\unsure{\mathrm{r}|\mathrm{n}}},\mathrm{E}_{\handwritten{6}},\mathrm{E}_{\handwritten{7}},\mathrm{E}_{\handwritten{8}}$ are those corresponding to Klein singularities), and $\handwritten{<}\mathrm{r}-1$ in \inserted{\handwritten{all}} other cases; in the first case, I would expect that there is just one orbit as stated, by the way, and maybe this property will be caracteristic of the ``homogene\deleted{\unsure{a|u}}ous'' case when all roots have same length. Also, what can be said of the dim\deleted{\unsure{\$|S}}ension of the set of all Borel subgroups containing a given element $\mathrm{g}\handwritten{\in}\mathrm{G}$, resp. whose Lie algebra contains a given element $\mathrm{x}\handwritten{\in}\underline{\mathrm{g}}$ ? If $2\mathrm{N}=\mathrm{n}-\mathrm{r}$ is the number of roots and if the dimension of the orbit of the given element is $2\mathrm{N}-2\mathrm{i}$, I would expect the dimension of the previous variety to be equal to $\mathrm{i}$; the inequality $\handwritten{\leqslant}\mathrm{i}$ would imply, by generalizing a depth argument of \unsure{Br\deleted{e}i\deleted{i}eskorn|Br\deleted{i}e\deleted{e}iskorn}, that the singularities of the fibers of your map $\underline{\mathrm{g}}\handwritten{\longrightarrow}\underline{\mathrm{h}}/\mathrm{W}$ (resp. Steinberg's $\mathrm{G}\handwritten{\longrightarrow}\mathrm{T}/\mathrm{W}$ in case $\mathrm{G}$ is simply connected instead of adjoint) are ``rational'', i.e. i\deleted{s}f $\mathrm{F}$ is such a fiber and $\mathrm{F}'\overset{\mathrm{f}}{\handwritten{\longrightarrow}}\mathrm{F}$ a resolution of singularities of $\mathrm{F}$, t\deleted{e}h\deleted{h}en $\mathrm{R}^{\mathrm{i}}\mathrm{f}_{\handwritten{*}}(\underline{\mathrm{O}}_{\mathrm{F}'}^{\deleted{\unsure{\mathrm{t}|*}}})=0$ for $\mathrm{i}\handwritten{>}0$. Maybe the ans\deleted{\illegible{1}}w\deleted{\illegible{1}}er \deleted{\unsure{ot|to}}\handwritten{to} most of the questions I am asking are well known and I am just ignorant; for some of them, it would indeed be scandalous that the experts do not know the answer ! In any case, I would be grateful for any comment you would have on any of my questions.

Sincerely yours\\
\handwritten{A\,\unsure{Grothendieck}}\\
% The vertical pen stroke at ``beinteresting'' is retained as a raised bar; no silent space has been supplied.
% archive sheet number: 22
% Isolated non-text vertical mark below the signature; no character reading assigned.
